\documentclass[12pt]{article}
\usepackage{amsmath,amssymb,amsthm}
\usepackage[margin=1in]{geometry}

\newtheorem{theorem}{Theorem}
\newtheorem{lemma}[theorem]{Lemma}

\title{Problem 285: Pythagorean Odds}
\author{Project Euler Solution}
\date{}

\begin{document}
\maketitle

\section{Problem Statement}

Albert picks $a,b$ uniformly from $[0,1]$, computes $\sqrt{(ka+1)^2+(kb+1)^2}$, and rounds to the nearest integer. He scores $k$ points if the result equals $k$. Find the expected total score for $k = 1$ to $100000$, rounded to 5 decimal places.

\section{Mathematical Foundation}

\begin{theorem}[Winning Probability as Annular Area]
The probability of winning round $k$ is $P_k = A(r_2,c) - A(r_1,c)$ where $r_1 = \frac{2k-1}{2k}$, $r_2 = \frac{2k+1}{2k}$, $c = \frac{1}{k}$, and
\[
A(r,c) = \begin{cases}
\displaystyle\frac{r^2}{2}\!\left(\frac{\pi}{2} - 2\arcsin\frac{c}{r}\right) - c\sqrt{r^2-c^2} + c^2 & \text{if } r \ge c\sqrt{2}, \\[8pt]
0 & \text{otherwise}.
\end{cases}
\]
\end{theorem}

\begin{proof}
The winning condition is equivalent to $r_1^2 \le x^2+y^2 \le r_2^2$ under the substitution $x = a + 1/k$, $y = b + 1/k$, where $(x,y) \in [c, 1+c]^2$.

Since $r_2 < 1+c$, the outer boundary of the annulus lies inside the square. The only binding constraints are $x \ge c$ and $y \ge c$. Define $A(r,c) = \mathrm{Area}\{(x,y): x \ge c,\, y \ge c,\, x^2+y^2 \le r^2\}$.

For $r \ge c\sqrt{2}$:
\[
A(r,c) = \int_c^{\sqrt{r^2-c^2}} \bigl(\sqrt{r^2-x^2} - c\bigr)\,dx.
\]

Using $\int \sqrt{r^2-x^2}\,dx = \frac{x\sqrt{r^2-x^2}}{2} + \frac{r^2}{2}\arcsin\frac{x}{r} + C$ and evaluating at the limits:
\begin{align*}
A(r,c) &= \left[\frac{x\sqrt{r^2-x^2}}{2} + \frac{r^2}{2}\arcsin\frac{x}{r} - cx\right]_c^{\sqrt{r^2-c^2}} \\
&= \frac{r^2}{2}\!\left(\arcsin\frac{\sqrt{r^2-c^2}}{r} - \arcsin\frac{c}{r}\right) - c\sqrt{r^2-c^2} + c^2.
\end{align*}

Applying $\arcsin\frac{\sqrt{r^2-c^2}}{r} = \frac{\pi}{2} - \arcsin\frac{c}{r}$ yields the stated formula.
\end{proof}

\section{Algorithm}

\begin{verbatim}
total = 0
for k = 1 to 100000:
    c = 1/k; r1 = (2k-1)/(2k); r2 = (2k+1)/(2k)
    P_k = area(r2, c) - area(r1, c)
    total += k * P_k
return round(total, 5)
\end{verbatim}

\section{Complexity Analysis}

\begin{itemize}
\item \textbf{Time:} $O(N)$ with $N = 100000$.
\item \textbf{Space:} $O(1)$.
\end{itemize}

\section{Answer}

\[
\boxed{157055.80999}
\]

\end{document}
